% Part: first-order-logic
% Chapter: sequent-calculus
% Section: equality

\documentclass[../../../include/open-logic-section]{subfiles}

\begin{document}

\olfileid{fol}{seq}{ide}

\olsection{\usetoken{P}{derivation} with \usetoken{S}{identity}}

包含等词的推导需要额外的初始矢列和推理规则。

\begin{defn}[$\eq$ 的初始矢列]
若 $t$ 是闭项，则 ${} \Sequent \eq[t][t]$ 是初始矢列。
\end{defn}

$\eq$ 的规则如下（$t_1$ 和 $t_2$ 是闭项）：

\begin{defish}
\Axiom$ \eq[t_1][t_2], \Gamma \fCenter \Delta, !A(t_1) $
\RightLabel{$\eq$}
\UnaryInf$\eq[t_1][t_2], \Gamma \fCenter \Delta, !A(t_2)$
\DisplayProof
\hfill
\Axiom$\eq[t_1][t_2], \Gamma \fCenter \Delta, !A(t_2) $
\RightLabel{$\eq$}
\UnaryInf$\eq[t_1][t_2], \Gamma  \fCenter \Delta, !A(t_1)$
\DisplayProof
\end{defish}

\begin{ex}
若 $s$ 和 $t$ 是闭项，则 $\eq[s][t], !A(s)
\Proves !A(t)$：
\begin{prooftree}
\Axiom$ !A(s) \fCenter !A(s)$
\RightLabel{\LeftR{\Weakening}}
\UnaryInf$\eq[s][t], !A(s)  \fCenter !A(s)$
\RightLabel{$\eq$}
\UnaryInf$\eq[s][t], !A(s)  \fCenter !A(t)$
\end{prooftree}
这或许就是大家所熟悉的同一物的可替换性原理，即 Leibniz（莱布尼茨） 律。

$\Log{LK}$ 可证明 $\eq$ 是对称且传递的：
\begin{prooftree}
\Axiom$ \fCenter \eq[t_1][t_1] $
\RightLabel{\LeftR{\Weakening}}
\UnaryInf$ \eq[t_1][t_2] \fCenter \eq[t_1][t_1] $
\RightLabel{$\eq$}
\UnaryInf$ \eq[t_1][t_2] \fCenter \eq[t_2][t_1]$
\DisplayProof\qquad\bottomAlignProof
\Axiom$ \eq[t_1][t_2] \fCenter \eq[t_1][t_2] $
\RightLabel{\LeftR{\Weakening}}
\UnaryInf$\eq[t_2][t_3], \eq[t_1][t_2]  \fCenter \eq[t_1][t_2] $
\RightLabel{$\eq$}
\UnaryInf$\eq[t_2][t_3], \eq[t_1][t_2]  \fCenter \eq[t_1][t_3]$
\RightLabel{\LeftR{\Exchange}}
\UnaryInf$\eq[t_1][t_2], \eq[t_2][t_3]  \fCenter \eq[t_1][t_3]$
\end{prooftree}
在左边的推导中，公式~$\eq[x][t_1]$ 就是我们的 $!A(x)$。在右边，我们取 $!A(x)$ 为~$\eq[t_1][x]$。
\end{ex}

\begin{prob}
给出下列矢列的推导：
\begin{enumerate}
\item $\Sequent \lforall[x][\lforall[y][((x = y \land !A(x)) \lif !A(y))]]$
\item $\lexists[x][!A(x)] \land \lforall[y][\lforall[z][((!A(y) \land
    !A(z)) \lif y = z)]] \Sequent 
\lexists[x][(!A(x) \land \lforall[y][(!A(y) \lif y = x)])]$
\end{enumerate}
\end{prob}

\end{document}